\ifajar
  \section{Latihan}
\else
  \section{Bab 4: Caching, Consistency, dan Resilience}
\fi

Seluruh file yang dipakai pada latihan ini tersedia di
\berkas{sources/chapter04/}, dihitung relatif terhadap akar repositori. Basis
data, Redis, dan layanan biaya kirim dijalankan di komputer sendiri, sehingga
tidak ada sistem milik pihak lain yang dibebani. Satu-satunya permintaan ke
sistem luar ada pada Latihan 1, yaitu tiga permintaan ke satu aset situs
kampus, dan ketiganya bersifat mengamati. Skrip yang sengaja membebani, yaitu
\texttt{stampede.py}, \texttt{resilience.py}, dan \texttt{ukur-ttfb.py}, tidak
boleh diarahkan ke sistem milik orang lain.

Angka pada tabel contoh diukur pada komputer uji dengan delapan inti prosesor.
Hasil tiap mahasiswa akan berbeda. Yang dilaporkan adalah pola perbandingannya
beserta penjelasannya, bukan angka persisnya.

Penyiapannya dilakukan sekali, seperti pada Listing \ref{lst:c4-siapkan}.
Skrip \texttt{siapkan-data.sh} memakai ulang skema, data uji, dan indeks
milik Bab 3.

\begin{lstlisting}[language=bash, caption={Menyiapkan basis data dan Redis, memakai \berkas{sources/chapter04/siapkan-data.sh}}, label={lst:c4-siapkan}]
cd sources/chapter04
docker compose up -d
./siapkan-data.sh

# Keluarannya (dipotong):
# === 02-data-uji.sql
#     tabel     | count
# --------------+--------
#  kategori     |     20
#  pelanggan    |  20000
#  buku         |  50000
#  pesanan      | 200000
#  item_pesanan | 600000
\end{lstlisting}

Bab ini menambahkan pustaka \texttt{redis} pada venv bersama di
\berkas{sources/.venv}, seperti pada Listing \ref{lst:c4-venv}.

\begin{lstlisting}[language=bash, caption={Menyiapkan venv untuk seluruh skrip Python}, label={lst:c4-venv}]
cd sources
python3 -m venv .venv
.venv/bin/pip install "psycopg[binary,pool]" sqlalchemy redis
\end{lstlisting}

\ifajar
  \subsection*{Latihan 1: Membaca Kebijakan Cache}
\else
  \subsection*{Prosedur 1: Membaca Kebijakan Cache}
\fi

Latihan ini membuktikan bahwa kebijakan \textit{cache} dapat dibaca langsung
dari header, dan bahwa validasi ulang menghemat isi tetapi tidak menghemat
perjalanan. Permintaan yang membawa \texttt{If-None-Match} disebut permintaan
bersyarat.

\begin{enumerate}
  \item Jalankan \texttt{python aplikasi.py} pada terminal pertama, lalu
        biarkan menyala sampai Latihan 2 selesai.
  \item Baca header halaman dan aset aplikasi, lalu catat
        \texttt{Cache-Control} dan \texttt{ETag}-nya. Perintah untuk halaman
        ada pada Listing \ref{lst:curl-etag}, dan perintah untuk aset ada pada
        Listing \ref{lst:c4-lat1-jalan}.
  \item Kirim permintaan bersyarat ke halaman dengan \texttt{ETag} dari langkah
        kedua, lalu catat status dan jumlah bytenya.
  \item Kirim satu permintaan \texttt{HEAD} ke aset \texttt{stylenew.css}
        milik situs kampus. Catat \texttt{date} dan \texttt{last-modified},
        lalu hitung masa \textit{fresh} tebakannya, yaitu sepersepuluh selang di antara
        keduanya.
  \item Kirim satu permintaan bersyarat ke aset kampus yang sama, memakai
        \texttt{ETag} dari langkah keempat.
  \item Unduh halaman aplikasi dan aset kampus tanpa header bersyarat, lalu
        catat status beserta jumlah byte yang diterima. Perintahnya ada pada
        Listing \ref{lst:c4-lat1-jalan}.
  \item Buka \texttt{http://localhost:8000/} di browser dengan panel Network
        terbuka dan \emph{Disable cache} tidak dicentang. Muat ulang halaman
        dua kali, lalu catat kolom \emph{Status} dan \emph{Size} untuk halaman
        dan untuk \texttt{app.3f9a2c.js}.
\ifajar
  \item Isi \tabisi{tab:c4-lat1-hasil}. Tabel \ref{tab:c4-lat1-contoh}
        memperlihatkan contoh yang sudah terisi.
\else
  \item Bandingkan hasil pengukuran sendiri dengan Tabel
        \ref{tab:c4-lat1-contoh}.
\fi
  \item Jawab tiga pertanyaan berikut. Pertama, mengapa halaman selalu
        ditanyakan ulang, sedangkan aset aplikasi tidak pernah ditanyakan
        ulang? Kedua, berapa byte yang dihemat oleh \textit{response} 304 pada halaman
        dan pada aset kampus, dan berapa RTT yang dihemat? Ketiga, apa yang
        tertulis pada kolom \emph{Size} untuk \texttt{app.3f9a2c.js} pada muat
        ulang kedua, dan mengapa?
\end{enumerate}

\begin{lstlisting}[language=bash, caption={Membaca header dan mengunduh isi, dari aplikasi sendiri dan situs kampus}, label={lst:c4-lat1-jalan}]
curl -sI http://localhost:8000/aset/app.3f9a2c.js \
  | grep -iE "^(http|cache-control)"
curl -sI https://www.pradita.ac.id/assets/front/css/stylenew.css \
  | grep -iE "^(date|last-modified|etag)"
curl -s -o /dev/null -w "%{http_code} %{size_download} byte\n" \
  -H 'If-None-Match: "67d3dca2-965"' \
  https://www.pradita.ac.id/assets/front/css/stylenew.css
curl -s -o /dev/null -w "%{http_code} %{size_download} byte\n" \
  http://localhost:8000/
curl -s -o /dev/null -w "%{http_code} %{size_download} byte\n" \
  https://www.pradita.ac.id/assets/front/css/stylenew.css

# Keluarannya:
# HTTP/1.0 200 OK
# Cache-Control: public, max-age=31536000, immutable
# date: Thu, 10 Sep 2026 08:24:59 GMT
# last-modified: Fri, 14 Mar 2025 07:37:06 GMT
# etag: "67d3dca2-965"
# 304 0 byte
# 200 331 byte
# 200 2405 byte
\end{lstlisting}

\begin{table}[!htb]
  \centering
  \caption{Contoh hasil Latihan 1 yang sudah terisi}
  \label{tab:c4-lat1-contoh}
  \begin{tabularx}{\textwidth}{lXlll}
    \hline
    \textbf{Sumber} & \textbf{\texttt{Cache-Control}} & \textbf{Masa \textit{fresh}} & \textbf{Biasa} & \textbf{Bersyarat} \\
    \hline
    Halaman aplikasi & \texttt{no-cache} & 0 detik & 200, 331 byte & 304, 0 byte \\
    \hline
    Aset aplikasi & \texttt{public, max-age=31536000, immutable} & 1 tahun & 200 & tanpa \texttt{ETag} \\
    \hline
    Aset kampus & tidak ada & sekitar 54 hari & 200, 2.405 byte & 304, 0 byte \\
    \hline
  \end{tabularx}
\end{table}

\ifajar
\begin{table}[!htb]
  \centering
  \caption{Lembar isian hasil Latihan 1}
  \label{tab:c4-lat1-hasil}
  \begin{tabularx}{\textwidth}{lXlll}
    \hline
    \textbf{Sumber} & \textbf{\texttt{Cache-Control}} & \textbf{Masa \textit{fresh}} & \textbf{Biasa} & \textbf{Bersyarat} \\
    \hline
    Halaman aplikasi & \dots & \dots & \dots & \dots \\
    \hline
    Aset aplikasi & \dots & \dots & \dots & \dots \\
    \hline
    Aset kampus & \dots & \dots & \dots & \dots \\
    \hline
    Aset lain pilihan sendiri & \dots & \dots & \dots & \dots \\
    \hline
  \end{tabularx}
\end{table}
\fi

Satu hal pada contoh tersebut layak diperhatikan. \texttt{ETag} aset kampus,
yaitu \texttt{67d3dca2-965}, tersusun dari dua bilangan heksadesimal. Bilangan
\texttt{965} sama dengan 2.405, yaitu ukuran file-nya dalam byte. Bilangan
\texttt{67d3dca2} sama dengan 1.741.937.826 detik sejak 1 Januari 1970, yaitu
waktu \texttt{last-modified} file tersebut. Server kampus menurunkan tanda
versinya dari waktu ubah dan ukuran file, tidak dari isinya seperti
\texttt{aplikasi.py}. Masa \textit{fresh} pada baris ketiga bergantung pada tanggal
pemeriksaan, sehingga angka tiap mahasiswa akan berbeda.

\ifajar
  \subsection*{Latihan 2: Cache-Aside dan Cache yang Macet}
\else
  \subsection*{Prosedur 2: Cache-Aside dan Cache yang Macet}
\fi

Latihan ini membuktikan bahwa \textit{cache-aside} memangkas waktu tunggu
server, dan bahwa \textit{cache} yang macet justru memperlambat aplikasi.

\begin{enumerate}
  \item Pastikan \texttt{aplikasi.py} masih menyala, lalu kosongkan kunci
        buku terlaris di Redis dengan perintah pertama pada Listing
        \ref{lst:c4-lat2-jalan}.
  \item Ukur \texttt{/api/terlaris} dan \texttt{/api/terlaris-cache}, masing-masing
        50 kali, dengan perintah pada listing yang sama.
  \item Bekukan Redis dengan \texttt{docker compose pause cache}, ukur
        \texttt{/api/terlaris-cache} 20 kali, lalu cairkan lagi dengan
        \texttt{docker compose unpause cache}.
\ifajar
  \item Isi \tabisi{tab:c4-lat2-hasil}. Tabel \ref{tab:c4-lat2-contoh}
        memperlihatkan contoh yang sudah terisi.
\else
  \item Bandingkan hasil pengukuran sendiri dengan Tabel
        \ref{tab:c4-lat2-contoh}.
\fi
  \item Jawab tiga pertanyaan berikut. Pertama, mengapa nilai terbesar pada
        baris \textit{cache-aside} hampir sama dengan p50 baris tanpa
        \textit{cache}? Kedua, dari mana angka sekitar 185 milidetik pada
        baris Redis dibekukan berasal? Ketiga, bandingkan p50 baris pertama
        dengan kolom tunggu server situs kampus pada Latihan 2 di Bab 2. Bagian
        mana dari waktu tunggu itu yang dapat dipangkas dengan cara pada
        latihan ini?
\end{enumerate}

\begin{lstlisting}[language=bash, caption={Mengukur waktu sampai byte pertama, tersedia di \berkas{sources/chapter04/ukur-ttfb.py}}, label={lst:c4-lat2-jalan}]
cd sources/chapter04
source ../.venv/bin/activate
docker compose exec -T cache redis-cli DEL terlaris:30hari
python ukur-ttfb.py http://localhost:8000/api/terlaris 50
python ukur-ttfb.py http://localhost:8000/api/terlaris-cache 50

# Keluarannya (dipotong):
# p50          :    70.44 ms
# p95          :    81.99 ms
# Terbesar     :    94.41 ms
# X-Cache      : - 50
# p50          :     0.31 ms
# p95          :     0.63 ms
# Terbesar     :    67.55 ms
# X-Cache      : MISS 1, HIT 49
\end{lstlisting}

\begin{table}[!htb]
  \centering
  \caption{Contoh hasil Latihan 2 yang sudah terisi, dalam milidetik}
  \label{tab:c4-lat2-contoh}
  \begin{tabularx}{\textwidth}{Xrrrl}
    \hline
    \textbf{Keadaan} & \textbf{p50} & \textbf{p95} & \textbf{Terbesar} & \textbf{X-Cache} \\
    \hline
    Tanpa \textit{cache} & 70,44 & 81,99 & 94,41 & tidak ada \\
    \hline
    \textit{Cache-aside} & 0,31 & 0,63 & 67,55 & MISS 1, HIT 49 \\
    \hline
    Redis dibekukan & 184,72 & 188,85 & 189,07 & MISS 20 \\
    \hline
  \end{tabularx}
\end{table}

\ifajar
\begin{table}[!htb]
  \centering
  \caption{Lembar isian hasil Latihan 2, dalam milidetik}
  \label{tab:c4-lat2-hasil}
  \begin{tabularx}{\textwidth}{Xrrrl}
    \hline
    \textbf{Keadaan} & \textbf{p50} & \textbf{p95} & \textbf{Terbesar} & \textbf{X-Cache} \\
    \hline
    Tanpa \textit{cache} & \dots & \dots & \dots & \dots \\
    \hline
    \textit{Cache-aside} & \dots & \dots & \dots & \dots \\
    \hline
    Redis dibekukan & \dots & \dots & \dots & \dots \\
    \hline
  \end{tabularx}
\end{table}
\fi

Baris ketiga pada contoh terisi tampak janggal, karena \textit{cache} yang
seharusnya mempercepat justru membuat \textit{response} dua setengah kali lebih lambat.
Penyebabnya dua kali penantian 50 milidetik, sekali untuk GET dan sekali untuk
SET, ditambah \textit{query} yang tetap harus dijalankan. Aplikasi tetap menjawab
seluruh 20 permintaan karena \textit{error} Redis ditelan dan \textit{retry} bawaan
pustakanya dimatikan.

\ifajar
  \subsection*{Latihan 3: Cache Stampede}
\else
  \subsection*{Prosedur 3: Cache Stampede}
\fi

Latihan ini membuktikan bahwa kunci yang kedaluwarsa dapat menghantam basis
data dengan banyak \textit{query} yang sama, lalu membandingkan dua cara meredamnya.

\begin{enumerate}
  \item Jalankan \texttt{stampede.py} tanpa argumen, seperti pada Listing
        \ref{lst:c4-lat3-jalan}. Skrip melepas 50 permintaan serentak.
  \item Catat jumlah \textit{query}, p50, p95, dan nilai terbesar untuk ketiga cara.
  \item Ulangi dengan 10 dan 100 permintaan serentak, yaitu
        \texttt{python stampede.py 10} dan \texttt{python stampede.py 100}.
  \item Jalankan sekali lagi dengan satu permintaan saja, yaitu
        \texttt{python stampede.py 1}, lalu bandingkan baris
        \texttt{stale\_while\_revalidate} dengan hasil 50 permintaan.
\ifajar
  \item Isi \tabisi{tab:c4-lat3-hasil}. Tabel \ref{tab:c4-lat3-contoh}
        memperlihatkan contoh yang sudah terisi.
\else
  \item Bandingkan hasil pengukuran sendiri dengan Tabel
        \ref{tab:c4-lat3-contoh}.
\fi
  \item Jawab tiga pertanyaan berikut. Pertama, mengapa p50
        \texttt{tanpa\_pengaman} jauh lebih besar daripada waktu satu \textit{query},
        padahal \textit{query}-nya sama persis? Hubungkan jawabannya dengan ukuran \textit{connection
        pool}. Kedua, mengapa jumlah \textit{query} pada \texttt{single\_flight}
        tetap satu berapa pun jumlah permintaannya? Ketiga, apa yang
        dikorbankan \texttt{stale\_while\_revalidate} untuk kecepatannya?
\end{enumerate}

\begin{lstlisting}[language=bash, caption={Memperagakan \textit{cache stampede}, tersedia di \berkas{sources/chapter04/stampede.py}}, label={lst:c4-lat3-jalan}]
cd sources/chapter04
source ../.venv/bin/activate
python stampede.py

# Keluarannya:
# 50 permintaan serentak tepat setelah kunci kedaluwarsa
#
#   Cara                    Query    p50 (ms)    p95 (ms)   Maks (ms)
#   tanpa_pengaman             50       895.7      1529.3      1564.3
#   single_flight               1       137.5       143.2       143.7
#   stale_while_revalidate      1         8.8        12.6        14.0
\end{lstlisting}

\begin{table}[!htb]
  \centering
  \caption{Contoh hasil Latihan 3 yang sudah terisi, waktu dalam milidetik}
  \label{tab:c4-lat3-contoh}
  \begin{tabularx}{\textwidth}{Xrrrr}
    \hline
    \textbf{Cara} & \textbf{\textit{Query}, 50} & \textbf{p50, 50} & \textbf{p95, 50} & \textbf{p50, 1} \\
    \hline
    \texttt{tanpa\_pengaman} & 50 & 895,7 & 1.529,3 & 135,2 \\
    \hline
    \texttt{single\_flight} & 1 & 137,5 & 143,2 & 125,1 \\
    \hline
    \texttt{stale\_while\_revalidate} & 1 & 8,8 & 12,6 & 0,6 \\
    \hline
  \end{tabularx}
\end{table}

\ifajar
\begin{table}[!htb]
  \centering
  \caption{Lembar isian hasil Latihan 3, waktu dalam milidetik}
  \label{tab:c4-lat3-hasil}
  \begin{tabularx}{\textwidth}{Xrrrrr}
    \hline
    \textbf{Cara} & \textbf{\textit{Query}, 100} & \textbf{p95, 10} & \textbf{p95, 50} & \textbf{p95, 100} & \textbf{p50, 1} \\
    \hline
    \texttt{tanpa\_pengaman} & \dots & \dots & \dots & \dots & \dots \\
    \hline
    \texttt{single\_flight} & \dots & \dots & \dots & \dots & \dots \\
    \hline
    \texttt{stale\_while\_revalidate} & \dots & \dots & \dots & \dots & \dots \\
    \hline
  \end{tabularx}
\end{table}
\fi

Angka pada judul kolom menyatakan jumlah permintaan serentak. Kolom terakhir
pada contoh terisi menjelaskan kejanggalan pada baris ketiga. Dengan satu
permintaan, nilai \textit{stale} terbaca dalam 0,6 milidetik. Dengan 50 permintaan, angka
itu naik menjadi 8,8 milidetik karena 50 utas program uji bergantian berjalan,
sedangkan Redis sendiri tetap cepat. Baris pertama paling bergantung pada
komputer uji, karena ditentukan oleh jumlah inti prosesor yang diperebutkan
\textit{query}.

\ifajar
  \subsection*{Latihan 4: Data Stale Setelah Perubahan}
\else
  \subsection*{Prosedur 4: Data Stale Setelah Perubahan}
\fi

Latihan ini mengukur berapa lama pengguna masih melihat harga lama setelah
harga diubah, pada empat cara memperbarui \textit{cache}.

\begin{enumerate}
  \item Jalankan \texttt{invalidasi.py}, seperti pada Listing
        \ref{lst:c4-lat4-jalan}.
  \item Catat jumlah pembacaan \textit{stale} dan lama \textit{stale}-nya untuk keempat cara.
  \item Ubah \texttt{TTL\_DETIK} di dalam \texttt{invalidasi.py} menjadi 10,
        jalankan ulang, lalu bandingkan dengan hasil pertama.
  \item Kembalikan \texttt{TTL\_DETIK} menjadi 5, ubah
        \texttt{JEDA\_HAPUS\_KEDUA\_DETIK} menjadi 1.0, jalankan ulang, lalu
        perhatikan baris terakhir.
\ifajar
  \item Isi \tabisi{tab:c4-lat4-hasil}. Tabel \ref{tab:c4-lat4-contoh}
        memperlihatkan contoh yang sudah terisi.
\else
  \item Bandingkan hasil pengukuran sendiri dengan Tabel
        \ref{tab:c4-lat4-contoh}.
\fi
  \item Jawab tiga pertanyaan berikut. Pertama, mengapa
        \texttt{hapus\_kalah\_race} sama buruknya dengan
        \texttt{ttl\_saja}? Kedua, baris mana yang berubah ketika TTL dinaikkan
        menjadi 10 detik, dan baris mana yang tidak berubah? Ketiga, bila
        pembaca paling lambat butuh 300 milidetik sejak membaca basis data
        sampai menyimpan ke \textit{cache}, berapa jeda penghapusan kedua yang
        aman?
\end{enumerate}

\begin{lstlisting}[language=bash, caption={Mengukur lama data stale, tersedia di \berkas{sources/chapter04/invalidasi.py}}, label={lst:c4-lat4-jalan}]
cd sources/chapter04
source ../.venv/bin/activate
python invalidasi.py

# Keluarannya:
# Harga 50000 menjadi 55000, TTL 5 detik, dibaca tiap 0.1 detik
#
#   Cara                  Baca stale Lama stale (s)
#   ttl_saja                      50            5.1
#   hapus_saat_menulis             0            0.0
#   hapus_kalah_race              50            5.1
#   hapus_dua_kali                 5            0.5
\end{lstlisting}

\begin{table}[!htb]
  \centering
  \caption{Contoh hasil Latihan 4 yang sudah terisi, TTL 5 detik}
  \label{tab:c4-lat4-contoh}
  \begin{tabularx}{\textwidth}{Xrr}
    \hline
    \textbf{Cara} & \textbf{Pembacaan \textit{stale}} & \textbf{Lama \textit{stale} (detik)} \\
    \hline
    \texttt{ttl\_saja} & 50 & 5,1 \\
    \hline
    \texttt{hapus\_saat\_menulis} & 0 & 0,0 \\
    \hline
    \texttt{hapus\_kalah\_race} & 50 & 5,1 \\
    \hline
    \texttt{hapus\_dua\_kali} & 5 & 0,5 \\
    \hline
  \end{tabularx}
\end{table}

\ifajar
\begin{table}[!htb]
  \centering
  \caption{Lembar isian hasil Latihan 4, lama \textit{stale} dalam detik}
  \label{tab:c4-lat4-hasil}
  \begin{tabularx}{\textwidth}{Xrrr}
    \hline
    \textbf{Cara} & \textbf{TTL 5} & \textbf{TTL 10} & \textbf{TTL 5, jeda kedua 1,0} \\
    \hline
    \texttt{ttl\_saja} & \dots & \dots & \dots \\
    \hline
    \texttt{hapus\_saat\_menulis} & \dots & \dots & \dots \\
    \hline
    \texttt{hapus\_kalah\_race} & \dots & \dots & \dots \\
    \hline
    \texttt{hapus\_dua\_kali} & \dots & \dots & \dots \\
    \hline
  \end{tabularx}
\end{table}
\fi

Pada contoh terisi, baris pertama dan ketiga sama persis, termasuk kelebihan
0,1 detik di atas TTL yang berasal dari selang pemeriksaan. Kesamaan itu
disengaja oleh skrip, karena \textit{race condition}-nya dipaksa selalu
terjadi. Pada sistem sungguhan, \textit{race condition} semacam itu jarang
terjadi, tetapi pada lalu lintas yang tinggi tetap akan terjadi sesekali.

\ifajar
  \subsection*{Latihan 5: Batas Waktu, Retry, dan Circuit Breaker}
\else
  \subsection*{Prosedur 5: Batas Waktu, Retry, dan Circuit Breaker}
\fi

Latihan ini membuktikan bahwa \textit{retry} menolong ketika gangguan hanya
sesekali, tetapi memperparah keadaan ketika layanan kewalahan, dan bahwa
\textit{circuit breaker} memotong beban tersebut.

\begin{enumerate}
  \item Jalankan \texttt{python layanan-lambat.py} pada terminal kedua, lalu
        biarkan menyala.
  \item Jalankan \texttt{resilience.py} pada terminal lain, seperti pada
        Listing \ref{lst:c4-lat5-jalan}. Kedua skenario dijalankan berurutan.
  \item Catat persentase berhasil, jumlah panggilan, dan durasi untuk keempat
        cara pada kedua skenario.
  \item Ubah \texttt{AMBANG\_GAGAL\_BERUNTUN} di dalam \texttt{resilience.py}
        menjadi 3, jalankan ulang, lalu catat jumlah panggilan
        \texttt{circuit\_breaker} pada skenario seluruhnya lambat.
  \item Bekukan Redis, jalankan \texttt{redis-macet.py} seperti pada Listing
        \ref{lst:redis-macet}, lalu cairkan Redis kembali.
\ifajar
  \item Isi \tabisi{tab:c4-lat5-hasil}. Tabel \ref{tab:c4-lat5-contoh}
        memperlihatkan contoh yang sudah terisi.
\else
  \item Bandingkan hasil pengukuran sendiri dengan Tabel
        \ref{tab:c4-lat5-contoh}.
\fi
  \item Jawab tiga pertanyaan berikut. Pertama, pada skenario seluruhnya
        lambat, mengapa \texttt{tanpa\_batas\_waktu} berhasil 100 persen tetapi
        tetap disebut cara terburuk? Kedua, berapa kali lipat beban yang
        diterima layanan karena \textit{retry}, dan mengapa \textit{circuit breaker}
        memanggil layanan 10 kali, bukan 5? Ketiga, mengapa
        \texttt{redis-macet.py} dengan pengaturan bawaan membutuhkan beberapa
        detik, padahal batas waktunya 50 milidetik?
\end{enumerate}

\begin{lstlisting}[language=bash, caption={Membandingkan empat cara memanggil layanan yang lambat, tersedia di \berkas{sources/chapter04/resilience.py}}, label={lst:c4-lat5-jalan}]
cd sources/chapter04
source ../.venv/bin/activate
python layanan-lambat.py        # terminal kedua, biarkan menyala
python resilience.py            # terminal lain

# Keluarannya (dipotong, skenario seluruhnya lambat):
#   Cara                Berhasil  p50 (ms)  p95 (ms)  Panggilan Durasi (s)
#   tanpa_batas_waktu       100%    2002.7    2004.1        100       20.0
#   batas_waktu               0%     302.0     304.0        100        3.0
#   batas_waktu_retry         0%     973.6    1032.8        300        9.9
#   circuit_breaker           0%       0.0     318.8         10        0.4
\end{lstlisting}

\begin{table}[!htb]
  \centering
  \caption{Contoh hasil Latihan 5 yang sudah terisi}
  \label{tab:c4-lat5-contoh}
  \begin{tabularx}{\textwidth}{Xrrrrr}
    \hline
     & \multicolumn{2}{c}{\textbf{Sesekali lambat}} & \multicolumn{3}{c}{\textbf{Seluruhnya lambat}} \\
    \hline
    \textbf{Cara} & \textbf{Berhasil} & \textbf{Panggilan} & \textbf{Berhasil} & \textbf{Panggilan} & \textbf{Durasi (s)} \\
    \hline
    \texttt{tanpa\_batas\_waktu} & 100\% & 100 & 100\% & 100 & 20,0 \\
    \hline
    \texttt{batas\_waktu} & 82\% & 100 & 0\% & 100 & 3,0 \\
    \hline
    \texttt{batas\_waktu\_retry} & 100\% & 123 & 0\% & 300 & 9,9 \\
    \hline
    \texttt{circuit\_breaker} & 100\% & 121 & 0\% & 10 & 0,4 \\
    \hline
  \end{tabularx}
\end{table}

\ifajar
\begin{table}[!htb]
  \centering
  \caption{Lembar isian hasil Latihan 5}
  \label{tab:c4-lat5-hasil}
  \begin{tabularx}{\textwidth}{Xrrrrr}
    \hline
     & \multicolumn{2}{c}{\textbf{Sesekali lambat}} & \multicolumn{3}{c}{\textbf{Seluruhnya lambat}} \\
    \hline
    \textbf{Cara} & \textbf{Berhasil} & \textbf{Panggilan} & \textbf{Berhasil} & \textbf{Panggilan} & \textbf{Durasi (s)} \\
    \hline
    \texttt{tanpa\_batas\_waktu} & \dots & \dots & \dots & \dots & \dots \\
    \hline
    \texttt{batas\_waktu} & \dots & \dots & \dots & \dots & \dots \\
    \hline
    \texttt{batas\_waktu\_retry} & \dots & \dots & \dots & \dots & \dots \\
    \hline
    \texttt{circuit\_breaker} & \dots & \dots & \dots & \dots & \dots \\
    \hline
    \texttt{circuit\_breaker}, ambang 3 & \dots & \dots & \dots & \dots & \dots \\
    \hline
  \end{tabularx}
\end{table}
\fi

Tiga angka pada contoh terisi menyimpang dari hitungan di atas kertas, dan
ketiganya sudah dijelaskan pada Seksi \ref{sec:resilience}. Batas waktu saja
berhasil 82 persen, bukan tepat 80 persen. \textit{Retry} berhasil 100 persen,
sedikit di atas hitungan 99,2 persen. Kedua selisih itu berasal dari undian
acak pada 100 permintaan. \textit{Circuit breaker} memanggil layanan 10 kali walaupun
ambangnya lima,
karena kesepuluh \textit{worker} sudah terlanjur memanggil sebelum kegagalan kelima
tercatat.

Setelah seluruh latihan selesai, hentikan \texttt{aplikasi.py} dan
\texttt{layanan-lambat.py} dengan Ctrl+C, lalu hentikan container dengan
\texttt{docker compose down -v}. Opsi \texttt{-v} ikut menghapus volume
datanya, sehingga penyiapan berikutnya dimulai dari keadaan bersih.
